\section{Introduction}\label{sec:intro}

Measuring the effect of monetary policy on the economy is one of the central problems of empirical
macroeconomics. The difficulty is familiar. Policy responds to the state of the economy, which in turn responds to
policy, and the econometrician sees a few hundred monthly observations of this feedback loop. The dominant
response has been to impose structure, through structural vector autoregressions \citep{sims1980} or local
projections \citep{jorda2005}, and with it linearity in the effect of policy.

A second approach, developed by \citet{angrist2011} and \citet{angrist2018}, borrows from microeconometric policy
evaluation. Each FOMC meeting is a unit that selects into a policy choice (ease, hold, or tighten), and the
effect of that choice is recovered by inverse probability weighting (IPW) on a propensity score for the
Committee's decision. The approach is semiparametric and allows effects to be nonlinear and asymmetric. Like all
selection-on-observables designs, however, it rests on two assumptions. \emph{Unconfoundedness} requires that,
given the observed covariates, the decision be unrelated to the future path the economy would take under each
alternative. \emph{Overlap} requires every decision to have positive probability at every value of the covariates.

This thesis starts from the second assumption. Some policy actions are simply not on the table in some states
of the economy. The Federal Reserve's own documents record which: before every meeting, Board staff circulate the
Bluebook, which lists the policy alternatives the Committee will consider. I use these \emph{policy menus},
extracted from the Bluebooks by text analysis, to make three contributions.
\begin{enumerate}[leftmargin=*,itemsep=2pt]
  \item \textbf{Structural non-overlap is observable.} In 133 scheduled meetings between August 1989 and
  January 2006 the FOMC never chose an action that was not on its Bluebook menu, and two thirds of the menus exclude
  either easing or tightening. The failure of overlap is not a feature of a fitted model. It is recorded in
  the Committee's briefing documents.
  \item \textbf{Menu-constrained propensity score and sharp bounds.} I estimate the decision model as an
  ordered probit renormalised over the menu, so that the propensity score is exactly zero for an action that is
  not on the menu. I then derive the sharp identified set for the average effect of a rate increase or decrease,
  exploiting the fact that in a menu that contains the status quo but not the action (or the reverse) half of
  the comparison is still identified. A one-parameter relaxation bounds how far effects in the non-overlap menus
  may differ from the effect in the overlap region, and yields a breakdown value for each estimate's sign.
  \item \textbf{Evidence.} Imposing the menus raises the log likelihood of the decision model by almost 19 points
  without any additional parameters. However, the effects that are point-identified on the overlap region imply that
  rate increases are followed by higher industrial production and lower unemployment, even when the propensity
  score conditions on the Greenbook forecasts that \citet{romer2004} use to purge policy of its predictable
  component. Worst-case bounds that make no assumption on the missing menus are wide and contain zero.
\end{enumerate}

The negative result is informative. It shows that correcting the overlap problem, which is real and
documented, does not deliver credible effects of monetary policy under selection on observables in this sample:
what remains is a failure of unconfoundedness, most plausibly because the Committee acts on information about the
future economy that neither macro aggregates, market expectations nor staff forecasts fully capture. The bounds
framework makes this transparent, because it separates what the data identify from what must be assumed.

\paragraph{Related literature.} The potential-outcomes approach to monetary policy follows \citet{angrist2011},
\citet{angrist2018} and \citet{rambachan2019}. The consequences of limited overlap for IPW are studied by
\citet{khan2010}, \citet{crump2009}, \citet{ma2020} and, with many covariates, \citet{damour2021}; I study the
case in which overlap fails exactly and the region of failure is observed. The bounds follow \citet{manski1990,manski2003},
and inference follows \citet{imbens2004}. The text data build on the FOMC text-analysis project of Montiel Olea and
coauthors and relate to \citet{hansen2018}. The use of staff forecasts to control for the Federal Reserve's
information set follows \citet{romer2004}.

\paragraph{Outline.} Section~\ref{sec:data} describes the FOMC's decision process and the construction of the
policy menus. Section~\ref{sec:framework} sets out the framework and derives the bounds. Section~\ref{sec:estimation}
covers estimation and inference, and Section~\ref{sec:results} reports results. Section~\ref{sec:conclusion} concludes.
